\chapter{Resultados}
\label{chap:resultados}

Segundo o Guia UECE 2026 (2.2.3.3, item e), os resultados são a apresentação objetiva e clara dos dados relevantes obtidos com a pesquisa -- sem interpretação ainda, apenas o que foi encontrado. Eles podem ser apresentados em forma de tabelas, quadros e gráficos, conforme o que se deseja destacar, e devem estar organizados de modo a responder diretamente a cada objetivo específico definido na Introdução.

Logo em seguida (ou nas mesmas seções, a critério do autor -- Guia 2.2.3.3, item f) vem a discussão: é o momento de comparar os resultados obtidos com os de outros pesquisadores já citados na Fundamentação Teórica/Trabalhos Relacionados, relacionando causa e efeito e retomando o problema ou a hipótese apresentados na Introdução. Resultados e discussão podem constituir um tópico único ou seções distintas, conforme a preferência do autor.

Organize este capítulo em uma seção por objetivo específico (ou por etapa do experimento/estudo), substituindo as seções de exemplo abaixo.

\section{Resultados do Experimento A}
\label{sec:resultados-do-experimento-a}

Apresente aqui os dados obtidos relativos ao primeiro objetivo específico (ou primeira etapa da coleta), preferencialmente com apoio de tabelas, quadros ou gráficos. Comente o que os dados mostram antes de compará-los com a literatura.

\section{Resultados do Experimento B}
\label{sec:resultados-do-experimento-b}

Apresente aqui os dados obtidos relativos ao segundo objetivo específico (ou segunda etapa da coleta), seguindo a mesma lógica da seção anterior: primeiro descreva o que foi encontrado, depois discuta à luz dos trabalhos relacionados citados anteriormente.
